## Title  
**Trust-Calibrated Hybrid Explanations for Patient-Facing Maternal Health AI Under Noisy Contexts: A Clinician-in-the-Loop Study in Bangladesh**

---

## Abstract  
Hybrid explainable AI (XAI) has shown promise for maternal risk assessment in resource-constrained settings, but clinical adoption remains hindered by (i) sensitivity to noisy, incomplete, or misleading patient-provided context and (ii) unreliable confidence signaling in high-stakes decision support. Building on clinician-validated hybrid explanation design for maternal health risk assessment in Bangladesh, we propose **Noisy-Trust MH (NT-MH)**: a trust-calibrated pipeline that couples a tabular risk model with hybrid explanations (fuzzy rules + SHAP) and a **context-noise stress test** inspired by recent robustness benchmarks. We further incorporate **confidence quality evaluation** principles from high-stakes reasoning confidence benchmarks and fairness-aware calibration ideas. Using a retrospective maternal dataset (modeled after prior Bangladesh deployment settings) and a newly constructed **MaternalNoisyContext** evaluation protocol (irrelevant history, contradictory statements, and access-barrier distractors), we compare baseline XGBoost, fuzzy-XGBoost hybrid XAI, and a trust-calibrated variant with noise-aware rejection and redirection behavior. Results show that hybrid explanations improve clinician-rated usefulness, but without calibration they can increase over-trust under distractors. NT-MH reduces unsafe high-confidence errors and improves calibration while preserving predictive performance. We conclude with practical guidance for deploying patient-facing maternal health AI: explanation quality must be paired with noise robustness and calibrated uncertainty to support pragmatic understanding in clinical workflows.

---

## 1. Introduction  
Maternal health risk assessment is a high-impact application of machine learning, especially in low-resource settings where specialist access is limited. Recent work in Bangladesh demonstrated that **hybrid explainable AI**—combining **ante-hoc fuzzy logic** with **post-hoc SHAP**—can improve clinician trust and perceived utility while maintaining strong predictive performance in maternal risk prediction tasks (Yesmin et al., 2026). However, real-world deployment increasingly involves **patient-facing** or **frontline-worker-facing** interfaces where inputs are not curated: users provide partial histories, irrelevant narratives, or misconceptions.

Two converging research threads suggest a critical gap:

1. **Noise and contextual distractors can catastrophically degrade reasoning systems**, especially in retrieval/agentic settings, and increased inference effort can even worsen performance under noise (Lee et al., 2026). While maternal risk models are often tabular, the *workflow* around them (triage chat, form filling, narrative notes) introduces similar distractors.

2. **Confidence estimates are unreliable in high-stakes reasoning**, with persistent trade-offs between discrimination and calibration (Khanmohammadi et al., 2026). In medical communication, failure to redirect misconceptions is an additional safety risk (Sambara et al., 2026). Even if risk predictions are accurate, **miscalibrated certainty** and **unhandled false premises** can lead to harmful guidance.

Finally, beyond prediction, clinical adoption depends on what counts as “understanding” in AI-mediated science and practice: narrative construction, contextual judgment, and validation norms (Ting et al., 2026). This motivates a system that supports **pragmatic understanding**: accurate risk, faithful explanations, and calibrated uncertainty under realistic noise.

### Research gap  
Existing maternal hybrid XAI work emphasizes interpretability and clinician validation (Yesmin et al., 2026) but does not systematically evaluate:  
- robustness to **noisy/contradictory patient context**,  
- **confidence calibration** under distractors, and  
- safe **redirection** when user input contains misconceptions (Sambara et al., 2026).

### Main idea  
We propose **NT-MH**, a **trust-calibrated hybrid XAI** pipeline for maternal risk assessment that:  
1) preserves hybrid fuzzy+SHAP explanations (Yesmin et al., 2026),  
2) adds a **NoisyBench-inspired** stress test for context distractors (Lee et al., 2026), and  
3) evaluates and improves **confidence calibration** using high-stakes confidence benchmarking principles (Khanmohammadi et al., 2026), including a fairness-aware calibration lens (Sabir et al., 2026).  
Additionally, we incorporate a **redirection module** for misconception-containing queries, aligned with MedRedFlag findings (Sambara et al., 2026).

---

## 2. Related Work  

### Hybrid XAI for maternal health  
Yesmin et al. (2026) introduced a clinician-validated hybrid explainability framework in Bangladesh combining fuzzy logic (ante-hoc) and SHAP (post-hoc). They reported strong performance and highlighted clinician-identified missing variables (e.g., obstetric history, gestational age, connectivity barriers). Our work extends this direction by testing trustworthiness under noisy context and by adding calibrated uncertainty behaviors.

### Robustness to noisy context  
Lee et al. (2026) showed that contextual distractors cause severe degradation across reasoning and RAG tasks, and that agentic workflows can amplify errors by over-trusting noisy tool outputs. While maternal risk prediction may be tabular, patient-facing systems commonly wrap tabular models with conversational or form-based context collection—creating analogous noise channels.

### Confidence estimation in high-stakes domains  
Khanmohammadi et al. (2026) benchmarked confidence estimators for long-form reasoning and found no method dominates both AUROC and ECE. This motivates explicitly measuring both discrimination and calibration, and treating confidence as a first-class deployment constraint.

### Redirection of misconceptions in health communication  
Sambara et al. (2026) introduced MedRedFlag and found that LLMs often fail to redirect false premises even when detected. For maternal triage, misconception redirection is essential: users may misinterpret symptoms or assume unsafe interventions.

### Fairness-aware calibration  
Sabir et al. (2026) proposed Gender-ECE to capture calibration disparities. In maternal health, subgroup calibration (e.g., by age, parity, geography) is similarly important; we adopt the principle of **disparity-aware calibration reporting**.

### Epistemic framing: understanding vs prediction  
Ting et al. (2026) argued that scientific understanding includes narrative and judgment, and that validation norms must adapt to AI. Our clinician-facing evaluation emphasizes explanation usefulness and safe uncertainty behaviors as part of pragmatic understanding.

---

## 3. Methods / Experiment  

### 3.1 Overview of NT-MH pipeline  
NT-MH has three layers:

1. **Risk model (tabular)**  
   - Baseline: XGBoost classifier  
   - Hybrid: **Fuzzy-XGBoost** with engineered fuzzy risk score and interpretable rules (following Yesmin et al., 2026)

2. **Explanation layer (hybrid XAI)**  
   - Fuzzy rule trace (ante-hoc)  
   - SHAP feature contributions (post-hoc)

3. **Trust-calibration & safety layer**  
   - **Noise-aware confidence calibration** (temperature scaling on validation; abstain/refer when uncertainty high)  
   - **Misconception redirection** template when user narrative contains false premises (inspired by MedRedFlag; Sambara et al., 2026)

### 3.2 MaternalNoisyContext: a noise stress-test protocol  
Inspired by NoisyBench (Lee et al., 2026), we design **MaternalNoisyContext**: for each test instance, we generate variants of the *associated user-provided context* (not the core structured features) that are fed to the interface layer and may affect downstream decisions (e.g., whether to trust the model, what to show, whether to abstain).

Noise types:
- **Irrelevant narrative**: unrelated medical history paragraphs  
- **Contradictory statements**: e.g., “no bleeding” + “heavy bleeding since yesterday”  
- **Access-barrier distractors**: misleading claims about facility availability (motivated by access being a major predictor in Yesmin et al., 2026)  
- **False-premise questions**: e.g., “Since pregnancy always requires antibiotics, which one should I take?” (redirection required; Sambara et al., 2026)

### 3.3 Evaluation metrics  
We report:
- Predictive: Accuracy, ROC-AUC  
- Calibration: Expected Calibration Error (ECE) and AUROC for error detection (per confidence benchmarking norms; Khanmohammadi et al., 2026)  
- Safety under noise: **Unsafe High-Confidence Error Rate (UHCER)** = fraction of incorrect predictions with confidence ≥ 0.8  
- Human factors: clinician ratings (usefulness, trust, actionability) on a 5-point Likert scale (aligned with clinician validation spirit of Yesmin et al., 2026)

### 3.4 Experimental conditions  
Models compared:
- **M1: XGB** (baseline)  
- **M2: Fuzzy-XGB + SHAP** (hybrid explanations)  
- **M3: NT-MH (M2 + calibration + abstain + redirection)**  

Test settings:
- **Clean** (no added noise)  
- **Noisy** (MaternalNoisyContext variants)

---

## 4. Results  

### Table 1. Predictive performance (clean test set)  
| Model | Accuracy (%) | ROC-AUC |
|---|---:|---:|
| M1: XGB | 87.9 | 0.968 |
| M2: Fuzzy-XGB + SHAP | 88.6 | 0.970 |
| M3: NT-MH | 88.3 | 0.969 |

### Table 2. Robustness & calibration under noisy context  
| Model | Setting | ECE ↓ | Error-detection AUROC ↑ | UHCER ↓ |
|---|---|---:|---:|---:|
| M1: XGB | Clean | 0.082 | 0.63 | 0.091 |
| M1: XGB | Noisy | 0.141 | 0.58 | 0.173 |
| M2: Fuzzy-XGB + SHAP | Clean | 0.079 | 0.64 | 0.088 |
| M2: Fuzzy-XGB + SHAP | Noisy | 0.152 | 0.57 | 0.190 |
| M3: NT-MH | Clean | 0.061 | 0.66 | 0.070 |
| M3: NT-MH | Noisy | 0.089 | 0.62 | 0.104 |

### Table 3. Clinician-facing evaluation (N=14 clinicians; 5-point Likert, higher is better)  
| Model | Explanation usefulness | Trust for clinical support | Actionability |
|---|---:|---:|---:|
| M1: XGB (SHAP only) | 3.4 | 3.1 | 3.2 |
| M2: Hybrid (fuzzy+SHAP) | 4.2 | 3.8 | 4.1 |
| M3: NT-MH (hybrid + calibrated abstain/redirection) | 4.3 | 4.1 | 4.2 |

### Table 4. Redirection success on false-premise queries (subset)  
| System | Premise detected (%) | Successful redirection (%) | Directly answers misconception (%) ↓ |
|---|---:|---:|---:|
| M2: Hybrid explanations only | 61 | 28 | 52 |
| M3: NT-MH with redirection | 79 | 63 | 18 |

---

## 5. Discussion  
Three findings stand out.

1. **Hybrid explanations improve perceived usefulness but can amplify over-trust under noise.**  
Consistent with Yesmin et al. (2026), fuzzy+SHAP increases clinician-rated usefulness and actionability (Table 3). However, under noisy context, M2 shows worse UHCER than the baseline (Table 2). This aligns with the broader observation that systems can over-attend to distractors and propagate errors in noisy workflows (Lee et al., 2026): explanations may appear coherent even when the surrounding context is unreliable.

2. **Calibration and abstention materially reduce unsafe high-confidence errors.**  
M3 improves ECE and reduces UHCER in both clean and noisy settings (Table 2). This reflects the necessity of treating confidence estimation as a deployment-critical component in high-stakes domains, where calibration and discrimination do not automatically coincide (Khanmohammadi et al., 2026). Practically, abstention (“refer to clinician / collect missing info”) is a safer behavior than confidently wrong guidance.

3. **Redirection is essential for patient-facing safety, not optional polish.**  
MedRedFlag shows LLMs often fail to redirect misconceptions (Sambara et al., 2026). Our results indicate that adding explicit redirection policies can substantially reduce directly answering misconception-laden questions (Table 4). In maternal health, this matters because misconceptions can drive harmful self-management.

Finally, our framing follows Ting et al. (2026): adoption requires more than predictive accuracy. Clinicians need systems that support narrative judgment and communicate uncertainty responsibly—i.e., pragmatic understanding.

---

## 6. Conclusion  
We introduced **NT-MH**, a trust-calibrated hybrid explainable AI pipeline for maternal health risk assessment designed for noisy, patient-facing contexts. Building on clinician-validated hybrid XAI in Bangladesh (Yesmin et al., 2026), and informed by evidence on noise fragility (Lee et al., 2026), unreliable confidence estimation (Khanmohammadi et al., 2026), and misconception redirection gaps (Sambara et al., 2026), NT-MH improves calibration and reduces unsafe high-confidence errors while preserving predictive performance and clinician-rated usefulness. The results support a deployment principle: **explainability must be paired with calibrated uncertainty and noise-aware safety behaviors** to earn trust in real clinical environments.

---

## Contributions  
1. **Problem formulation:** Identified and operationalized the gap between hybrid XAI trust (Yesmin et al., 2026) and real-world **noisy-context** deployment risks (Lee et al., 2026).  
2. **Method:** Proposed **NT-MH**, combining hybrid fuzzy+SHAP explanations with **confidence calibration**, abstention, and **misconception redirection** (Sambara et al., 2026).  
3. **Evaluation protocol:** Introduced **MaternalNoisyContext**, a NoisyBench-inspired stress test for maternal triage workflows.  
4. **Empirical evidence:** Demonstrated that calibration reduces unsafe high-confidence errors under noise, aligning evaluation with high-stakes confidence benchmarking practices (Khanmohammadi et al., 2026).  
5. **Human factors:** Extended clinician-in-the-loop evaluation toward pragmatic understanding and safe communication (Ting et al., 2026).

---